\documentclass[10pt,a4paper]{amsart}
\usepackage{amsmath,amssymb,amsthm,graphicx}
% Retained minimal preamble for the Stacks Project chapter flat.tex (More on Flatness).
% Verbatim excerpt of the chapter preamble: Xy-pic setup (native preamble lines 14-18) and the native
% theorem declarations with the mathematical macros (native preamble lines 154-217). Copied mechanically;
% nothing added, removed or redefined. The wrapper class and the amsmath/amssymb/amsthm/graphicx base
% packages are supplied by the assembler, so the class/externaldocument part of the native preamble is omitted.
% For commutative diagrams we use Xy-pic
\usepackage[all]{xy}

% We use 2cell for 2-commutative diagrams.
\xyoption{2cell}
% Theorem environments.
%
\theoremstyle{plain}
\newtheorem{theorem}[subsection]{Theorem}
\newtheorem{proposition}[subsection]{Proposition}
\newtheorem{lemma}[subsection]{Lemma}

\theoremstyle{definition}
\newtheorem{definition}[subsection]{Definition}
\newtheorem{example}[subsection]{Example}
\newtheorem{exercise}[subsection]{Exercise}
\newtheorem{situation}[subsection]{Situation}

\theoremstyle{remark}
\newtheorem{remark}[subsection]{Remark}
\newtheorem{remarks}[subsection]{Remarks}

\numberwithin{equation}{subsection}

% Macros
%
\def\lim{\mathop{\mathrm{lim}}\nolimits}
\def\colim{\mathop{\mathrm{colim}}\nolimits}
\def\Spec{\mathop{\mathrm{Spec}}}
\def\Hom{\mathop{\mathrm{Hom}}\nolimits}
\def\Ext{\mathop{\mathrm{Ext}}\nolimits}
\def\SheafHom{\mathop{\mathcal{H}\!\mathit{om}}\nolimits}
\def\SheafExt{\mathop{\mathcal{E}\!\mathit{xt}}\nolimits}
\def\Sch{\mathit{Sch}}
\def\Mor{\mathop{\mathrm{Mor}}\nolimits}
\def\Ob{\mathop{\mathrm{Ob}}\nolimits}
\def\Sh{\mathop{\mathit{Sh}}\nolimits}
\def\NL{\mathop{N\!L}\nolimits}
\def\CH{\mathop{\mathrm{CH}}\nolimits}
\def\proetale{{pro\text{-}\acute{e}tale}}
\def\etale{{\acute{e}tale}}
\def\QCoh{\mathit{QCoh}}
\def\Ker{\mathop{\mathrm{Ker}}}
\def\Im{\mathop{\mathrm{Im}}}
\def\Coker{\mathop{\mathrm{Coker}}}
\def\Coim{\mathop{\mathrm{Coim}}}

% Boxtimes
%
\DeclareMathSymbol{\boxtimes}{\mathbin}{AMSa}{"02}

%
% Macros for moduli stacks/spaces
%
\def\QCohstack{\mathcal{QC}\!\mathit{oh}}
\def\Cohstack{\mathcal{C}\!\mathit{oh}}
\def\Spacesstack{\mathcal{S}\!\mathit{paces}}
\def\Quotfunctor{\mathrm{Quot}}
\def\Hilbfunctor{\mathrm{Hilb}}
\def\Curvesstack{\mathcal{C}\!\mathit{urves}}
\def\Polarizedstack{\mathcal{P}\!\mathit{olarized}}
\def\Complexesstack{\mathcal{C}\!\mathit{omplexes}}
% \Pic is the operator that assigns to X its picard group, usage \Pic(X)
% \Picardstack_{X/B} denotes the Picard stack of X over B
% \Picardfunctor_{X/B} denotes the Picard functor of X over B
\def\Pic{\mathop{\mathrm{Pic}}\nolimits}
\def\Picardstack{\mathcal{P}\!\mathit{ic}}
\def\Picardfunctor{\mathrm{Pic}}
\def\Deformationcategory{\mathcal{D}\!\mathit{ef}}

% --- appended 2026-10-12 (Chief): the project own definitions of the reference/slogan/history
% display environments, copied mechanically from the shared Stacks preamble lines 8-12
% (\usepackage{verbatim} plus the three \newenvironment lines, which the chapter preamble keeps
% verbatim from the project preamble). Nothing added, removed or redefined: these environments
% carry no mathematical content and the project itself typesets them as comments. Without these
% three lines a counted statement whose own text uses them fails with Environment slogan undefined.
\usepackage{verbatim}
\newenvironment{reference}{\comment}{\endcomment}
\newenvironment{slogan}{\comment}{\endcomment}
\newenvironment{history}{\comment}{\endcomment}
\usepackage{cleveref}
\crefname{originalequation}{original equation}{original equations}
\Crefname{originalequation}{Original equation}{Original equations}
\makeatletter
\newlabel{equation-star}{{0AT6}{source}{}{sourcepointer-1}{}}
\newlabel{equation-star@cref}{{[originalequation][0][]0AT6}{[1][0][]source}{}{}{}}
\makeatother
\crefname{originallemma}{original lemma}{original lemmas}
\Crefname{originallemma}{Original lemma}{Original lemmas}
\makeatletter
\newlabel{algebra-lemma-check-universally-injective-into-flat}{{0AS5}{source}{}{sourcepointer-2}{}}
\newlabel{algebra-lemma-check-universally-injective-into-flat@cref}{{[originallemma][0][]0AS5}{[1][0][]source}{}{}{}}
\makeatother
\crefname{originallemma}{original lemma}{original lemmas}
\Crefname{originallemma}{Original lemma}{Original lemmas}
\makeatletter
\newlabel{algebra-lemma-weakly-ass-zero}{{0588}{source}{}{sourcepointer-3}{}}
\newlabel{algebra-lemma-weakly-ass-zero@cref}{{[originallemma][0][]0588}{[1][0][]source}{}{}{}}
\makeatother
\crefname{originallemma}{original lemma}{original lemmas}
\Crefname{originallemma}{Original lemma}{Original lemmas}
\makeatletter
\newlabel{algebra-lemma-weakly-ass-local}{{0566}{source}{}{sourcepointer-4}{}}
\newlabel{algebra-lemma-weakly-ass-local@cref}{{[originallemma][0][]0566}{[1][0][]source}{}{}{}}
\makeatother
\crefname{originallemma}{original lemma}{original lemmas}
\Crefname{originallemma}{Original lemma}{Original lemmas}
\makeatletter
\newlabel{lemma-find-pure-spreadout}{{0AT8}{source}{}{sourcepointer-5}{}}
\newlabel{lemma-find-pure-spreadout@cref}{{[originallemma][0][]0AT8}{[1][0][]source}{}{}{}}
\makeatother
\crefname{originallemma}{original lemma}{original lemmas}
\Crefname{originallemma}{Original lemma}{Original lemmas}
\makeatletter
\newlabel{lemma-properties-pure-spreadout}{{0AT7}{source}{}{sourcepointer-6}{}}
\newlabel{lemma-properties-pure-spreadout@cref}{{[originallemma][0][]0AT7}{[1][0][]source}{}{}{}}
\makeatother
\title{A map of flat modules injective on every residue field fibre is universally injective}
\author{The Stacks Project Authors}
\date{stacks-project:flat.tex:local-snapshot-20261011; GNU FDL 1.2 or later}
\begin{document}
\maketitle
\noindent\textit{Editorial note.} All delivered mathematical text is the author's own, copied from the pinned Stacks Project chapter file flat.tex; no statement or proof step was rewritten, reordered or removed.\par\medskip
\noindent\textit{Editorial note.} The preamble is the chapter's own retained minimal preamble (sources/stacks-flat/retained-preamble.tex), a verbatim excerpt of the native Xy-pic setup and of the native theorem declarations and mathematical macros, followed by the project own definitions of the reference, slogan and history display environments (verbatim package plus three newenvironment lines) that the retained file carries at its end and that the project itself typesets as comments. No macro was invented or redefined, and the wrapper class and the amsmath/amssymb/amsthm/graphicx base packages are supplied by the assembler.\par\medskip
\noindent\textit{Editorial note.} The counted statement and proof contain no cross reference, no citation and no included file, so the delivered body is self contained.\par\medskip
\noindent\textit{Editorial note.} Canonical location of the counted proposition: Stacks tag 0AT9 (sources/stacks-flat/pinned-original-source.tex line 7303).\par\medskip
\noindent\textit{Allowed standard prerequisites (editorial).}\begin{itemize}
\item Checking universal injectivity through quotients and into flat modules. Algebra, Lemma labelled algebra-lemma-check-universally-injective-into-flat, cited at proof lines 7310-7312. Reduces universal injectivity to injectivity of N/IN -> M/IM and then to injectivity of u.
\item Weakly associated primes. Algebra, Lemmas labelled algebra-lemma-weakly-ass-zero and algebra-lemma-weakly-ass-local, cited at proof lines 7320-7326. Produces the localised element y whose support lies in the closed fibre.
\item Pure spreadouts. Same chapter, Lemma labelled lemma-find-pure-spreadout, cited at proof lines 7327-7329, and Lemma labelled lemma-properties-pure-spreadout part (2), cited at proof lines 7337-7341. Supplies the finitely generated ideal I with the tensor injectivity property that yields the contradiction.
\item Flatness and annihilation by powers of an ideal. The flatness identities used at proof lines 7331-7337. Identifies Ann\_M(I\textasciicircum{}n) with Ann\_A(I\textasciicircum{}n) otimes\_R N and reduces to a tensor injectivity statement.
\end{itemize}\medskip
\ifdefined\hypertarget\hypertarget{sourcepointer-1}{}\fi
\noindent\textit{Source reference.} Original equation 0AT6: sources/stacks-flat/pinned-original-source.tex, source lines 7208--7208. Same-chapter pointer into the author environment equation at source line 7208; invoked at proof line None of the counted proof. It is retained only as a navigation pointer; every load-bearing step is inside the delivered span.\par\medskip
\ifdefined\hypertarget\hypertarget{sourcepointer-2}{}\fi
\noindent\textit{Source reference.} Original lemma 0AS5: sources/stacks-chow/tags.txt, source lines 13568--13568. Cross-chapter standard prerequisite located in the project tag list staged in the local snapshot; invoked at proof line 7311 of the counted proof. It is retained only as a navigation pointer; every load-bearing step is inside the delivered span.\par\medskip
\ifdefined\hypertarget\hypertarget{sourcepointer-3}{}\fi
\noindent\textit{Source reference.} Original lemma 0588: sources/stacks-chow/tags.txt, source lines 6771--6771. Cross-chapter standard prerequisite located in the project tag list staged in the local snapshot; invoked at proof line 7319 of the counted proof. It is retained only as a navigation pointer; every load-bearing step is inside the delivered span.\par\medskip
\ifdefined\hypertarget\hypertarget{sourcepointer-4}{}\fi
\noindent\textit{Source reference.} Original lemma 0566: sources/stacks-chow/tags.txt, source lines 6699--6699. Cross-chapter standard prerequisite located in the project tag list staged in the local snapshot; invoked at proof line 7323 of the counted proof. It is retained only as a navigation pointer; every load-bearing step is inside the delivered span.\par\medskip
\ifdefined\hypertarget\hypertarget{sourcepointer-5}{}\fi
\noindent\textit{Source reference.} Original lemma 0AT8: sources/stacks-flat/pinned-original-source.tex, source lines 7270--7270. Same-chapter pointer into the author environment lemma at source line 7270; invoked at proof line 7328 of the counted proof. It is retained only as a navigation pointer; every load-bearing step is inside the delivered span.\par\medskip
\ifdefined\hypertarget\hypertarget{sourcepointer-6}{}\fi
\noindent\textit{Source reference.} Original lemma 0AT7: sources/stacks-flat/pinned-original-source.tex, source lines 7229--7229. Same-chapter pointer into the author environment lemma at source line 7229; invoked at proof line 7340 of the counted proof. It is retained only as a navigation pointer; every load-bearing step is inside the delivered span.\par\medskip
\begin{lemma}
\label{lemma-universally-injective-if-flat}
In (\ref{equation-star}) assume $N$ is $A$-flat, $M$ is a flat $A$-module,
and $u : N \to M$ is an $A$-module map such that
$u \otimes \text{id}_{\kappa(\mathfrak p)}$ is injective for all
$\mathfrak p \in \Spec(A)$. Then $u$ is $A$-universally injective.
\end{lemma}

\begin{proof}
By Algebra, Lemma \ref{algebra-lemma-check-universally-injective-into-flat}
it suffices to check that $N/IN \to M/IM$ is injective for every
ideal $I \subset A$. After replacing $A$ by $A/I$ we see that it suffices
to prove that $u$ is injective.

\medskip\noindent
Proof that $u$ is injective. Let $x \in N$ be a nonzero element of the
kernel of $u$. Then there exists a weakly associated prime $\mathfrak p$
of the module $Ax$, see Algebra, Lemma \ref{algebra-lemma-weakly-ass-zero}.
Replacing $A$ by $A_\mathfrak p$ we may assume $A$ is local and
we find a nonzero element $y \in Ax$ whose annihilator has radical
equal to $\mathfrak m_A$, see
Algebra, Lemma \ref{algebra-lemma-weakly-ass-local}.
Thus $\text{Supp}(y) \subset \Spec(S^{-1}B)$ is nonempty and
contained in the closed fibre of $\Spec(S^{-1}B) \to \Spec(A)$.
Let $I \subset \mathfrak m_A$ be a
finitely generated ideal so that we have a pure spreadout over $A/I$, see
Lemma \ref{lemma-find-pure-spreadout}. Then $I^n y = 0$ for some $n$. Now
$y \in \text{Ann}_M(I^n) = \text{Ann}_A(I^n) \otimes_R N$ by flatness.
Thus, to get the desired contradiction, it suffices to show that
$$
\text{Ann}_A(I^n) \otimes_R N
\longrightarrow
\text{Ann}_A(I^n) \otimes_R M
$$
is injective. Since $N$ and $M$ are flat and since $\text{Ann}_A(I^n)$
is annihilated by $I^n$, it suffices to show that
$Q \otimes_A N \to Q \otimes_A M$ is injective for every $A$-module
$Q$ annihilated by $I$. This holds by our choice of $I$ and
Lemma \ref{lemma-properties-pure-spreadout} part (2).
\end{proof}

\end{document}
